% Lösungen Einheit 2 von 2 – Linearkombination mit Nebenbedingung (zusammenbau v0.1)
\einheitenkopf{Einheit 2 von 2 · Linearkombination mit Nebenbedingung}

\erg{8}{a) Summe 1, beide in [0; 1] \quad b) Mit $p = 1 - q$: $\vec{OX} = \vec{OA} + q \cdot (\vec{OB} - \vec{OA}) = (1 | 0 | 2) + q \cdot (2 | 4 | -2)$. \quad c) $q = 0$: $X = A$; $q = 1$: $X = B$; $q = 0{,}5$: der Mittelpunkt $(3 | 1 | 3)$ der Strecke $AB$. \quad d) Mit $p = 1 - q$: $\vec{OX} = \vec{OB} + q \cdot (\vec{OS} - \vec{OB}) = \vec{OB} + q \cdot \vec{BS}$ mit $\vec{BS} = (-3 | 3 | 4)$. Das ist die Parameterform der Geraden $BS$; für $q \in [0; 1]$ läuft $X$ von $B$ ($q = 0$) bis $S$ ($q = 1$), also liegt $X$ auf der Kante $BS$.}
\erg{9}{a) Ein Zahlenbeispiel beweist nichts. Allgemein mit $p = 1 - q$: $\vec{OX} = \vec{OA} + q \cdot \vec{AB}$ mit $\vec{AB} = (2 | 4 | 2)$ und $q \in [0; 1]$: das ist die Strecke $AB$. \quad b) Einsetzen und Ausklammern ergibt $\vec{OA} + q \cdot (\vec{OB} - \vec{OA})$: fester Stützvektor $\vec{OA}$ plus Parameter mal Richtungsvektor $\vec{AB}$.}
